% GENERATED FILE. Edit canonical lessons, then run npm run generate:books.
% canonical-source-hash: fnv1a64:a8feff604dc96cc5
% canonical-lessons: GU-C451-valagvu, GU-C451-cumvu, GU-C451-bajhvu, GU-C451-pyar, GU-C451-capvu

\chapter{Holding Close}
\label{ch:gu-holding-close}
\hlchaptermodality{451}

\begin{chapteropening}
\textbf{By the end of this chapter:} \emph{I can say how a child holds on to its mother.}

Complete the last lesson of chapter 451: I can say how a child holds on to its mother.
\end{chapteropening}

\section[\texorpdfstring{\gu{વળગવું} (vaḷagvũ)}{વળગવું (vaḷagvũ)}]{\gu{વળગવું} (vaḷagvũ) — to cling to; to hug tight}

\label{lesson:GU-C451-valagvu}



\begin{quote}
Before the new one: say the Gujarati for to captivate, then the Gujarati for pampering.
\end{quote}

\subsection*{You'll want to know: \gu{વળગવું}}
\textbf{\gu{વળગવું}} — \emph{vaḷagvũ} — \textquotedblleft{}to cling to; to hug tight\textquotedblright{}.

\begin{quote}
\gu{બાળક} \gu{બીકથી} \gu{માતાને} \gu{વળગી} \gu{પડ્યું}.
\end{quote}

\begin{quote}
\emph{bāḷak bīkthī mātāne vaḷgī paḍyũ.}
\end{quote}

\begin{quote}
The child clung to its mother in fear.
\end{quote}

Past (he …): \gu{તે} \gu{માતાને} \gu{વળગી} \gu{પડ્યો}, \emph{te mātāne vaḷgī paḍyo}.

\begin{cousinweb}
\textbf{\gu{વળગવું}} (\emph{vaḷagvũ}) goes back to a reconstructed \emph{*vilagyati}, \textquotedblleft{}hangs on to\textquotedblright{}; CDIAL lists the Gujarati verb, \textquotedblleft{}to adhere to, seize\textquotedblright{} (CDIAL 11881).
\end{cousinweb}

\begin{grammarlens}[title={What you've built}]
One more word for describing what you lived and felt.
\end{grammarlens}

\subsection*{Your turn}
\begin{itemize}
\raggedright
  \item \emph{Say it:} \emph{vaḷagvũ}
  \item \emph{Say it:} \emph{vaḷagvũ}, once more
  \item \emph{Say it:} \emph{bāḷak bīkthī mātāne vaḷgī paḍyũ.} — The child clung to its mother in fear.
  \item \emph{Recall:} say the Gujarati for to captivate, then the Gujarati for pampering, then say \emph{vaḷagvũ} again
\end{itemize}

\begin{tcolorbox}[breakable,colback=teal!4,colframe=teal!35!black,title={Before you move on}]
What does \gu{વળગવું} mean? (\textquotedblleft{}To cling to; to hug tight\textquotedblright{}.) Say it once more.
\end{tcolorbox}
\section[\texorpdfstring{\gu{ચૂમવું} (cūmvũ)}{ચૂમવું (cūmvũ)}]{\gu{ચૂમવું} (cūmvũ) — to kiss}

\label{lesson:GU-C451-cumvu}



\begin{quote}
Before the new one: say the Gujarati for pampering, then the Gujarati for to cling.
\end{quote}

\subsection*{You'll want to know: \gu{ચૂમવું}}
\textbf{\gu{ચૂમવું}} — \emph{cūmvũ} — \textquotedblleft{}to kiss\textquotedblright{}.

\begin{quote}
\gu{માતા} \gu{રોજ} \gu{રાતે} \gu{બાળકના} \gu{કપાળને} \gu{ચૂમે} \gu{છે}.
\end{quote}

\begin{quote}
\emph{mātā roj rāte bāḷaknā kapāḷne cūme che.}
\end{quote}

\begin{quote}
Every night the mother kisses the child's forehead.
\end{quote}

Past, with \gu{તેણે} (he or she …): \gu{તેણે} \gu{બાળકનું} \gu{કપાળ} \gu{ચૂમ્યું}, \emph{teṇe bāḷaknũ kapāḷ cūmyũ}.

\begin{cousinweb}
\textbf{\gu{ચૂમવું}} (\emph{cūmvũ}) is inherited from Sanskrit \emph{cumbati}, \textquotedblleft{}kisses\textquotedblright{}; CDIAL lists the Gujarati verb (CDIAL 4870).
\end{cousinweb}

\begin{grammarlens}[title={What you've built}]
One more word for describing what you lived and felt.
\end{grammarlens}

\subsection*{Your turn}
\begin{itemize}
\raggedright
  \item \emph{Say it:} \emph{cūmvũ}
  \item \emph{Say it:} \emph{cūmvũ}, once more
  \item \emph{Say it:} \emph{mātā roj rāte bāḷaknā kapāḷne cūme che.} — Every night the mother kisses the child's forehead.
  \item \emph{Recall:} say the Gujarati for pampering, then the Gujarati for to cling, then say \emph{cūmvũ} again
\end{itemize}

\begin{tcolorbox}[breakable,colback=teal!4,colframe=teal!35!black,title={Before you move on}]
What does \gu{ચૂમવું} mean? (\textquotedblleft{}To kiss\textquotedblright{}.) Say it once more.
\end{tcolorbox}
\section[\texorpdfstring{\gu{બાઝવું} (bājhvũ)}{બાઝવું (bājhvũ)}]{\gu{બાઝવું} (bājhvũ) — to hold on tight; to grapple}

\label{lesson:GU-C451-bajhvu}



\begin{quote}
Before the new one: say the Gujarati for to cling, then the Gujarati for to kiss.
\end{quote}

\subsection*{You'll want to know: \gu{બાઝવું}}
\textbf{\gu{બાઝવું}} — \emph{bājhvũ} — \textquotedblleft{}to hold on tight; to grapple\textquotedblright{}.

\begin{quote}
\gu{નાની} \gu{દીકરી} \gu{પિતાના} \gu{પગને} \gu{બાઝી} \gu{પડી}.
\end{quote}

\begin{quote}
\emph{nānī dīkrī pitānā pagne bājhī paḍī.}
\end{quote}

\begin{quote}
The little daughter held on tight to her father's leg.
\end{quote}

Past (she …): \gu{તે} \gu{પિતાને} \gu{બાઝી} \gu{પડી}, \emph{te pitāne bājhī paḍī}.

\begin{cousinweb}
\textbf{\gu{બાઝવું}} (\emph{bājhvũ}) is inherited from Sanskrit \emph{badhyatē}, \textquotedblleft{}is bound\textquotedblright{}; CDIAL lists the Gujarati verb as \textquotedblleft{}to stick close to, catch hold of, fight\textquotedblright{} (CDIAL 9134).
\end{cousinweb}

\begin{grammarlens}[title={What you've built}]
One more word for describing what you lived and felt.
\end{grammarlens}

\subsection*{Your turn}
\begin{itemize}
\raggedright
  \item \emph{Say it:} \emph{bājhvũ}
  \item \emph{Say it:} \emph{bājhvũ}, once more
  \item \emph{Say it:} \emph{nānī dīkrī pitānā pagne bājhī paḍī.} — The little daughter held on tight to her father's leg.
  \item \emph{Recall:} say the Gujarati for to cling, then the Gujarati for to kiss, then say \emph{bājhvũ} again
\end{itemize}

\begin{tcolorbox}[breakable,colback=teal!4,colframe=teal!35!black,title={Before you move on}]
What does \gu{બાઝવું} mean? (\textquotedblleft{}To hold on tight; to grapple\textquotedblright{}.) Say it once more.
\end{tcolorbox}
\section[\texorpdfstring{\gu{પ્યાર} (pyār)}{પ્યાર (pyār)}]{\gu{પ્યાર} (pyār) — love; loving fondness}

\label{lesson:GU-C451-pyar}



\begin{quote}
Before the new one: say the Gujarati for to kiss, then the Gujarati for to hold on tight.
\end{quote}

\subsection*{You'll want to know: \gu{પ્યાર}}
\textbf{\gu{પ્યાર}} — \emph{pyār} — \textquotedblleft{}love; loving fondness\textquotedblright{}.

\begin{quote}
\gu{માતાના} \gu{પ્યારથી} \gu{બાળક} \gu{શાંત} \gu{થઈ} \gu{ગયું}.
\end{quote}

\begin{quote}
\emph{mātānā pyārthī bāḷak śānt thaī gayũ.}
\end{quote}

\begin{quote}
In its mother's love the child grew calm.
\end{quote}

A masculine noun.

\begin{cousinweb}
\textbf{\gu{પ્યાર}} (\emph{pyār}) goes back to Sanskrit \emph{priyakāra}, \textquotedblleft{}doing a kindness\textquotedblright{}; CDIAL lists the Gujarati word, \textquotedblleft{}love\textquotedblright{} (CDIAL 8975).
\end{cousinweb}

\begin{grammarlens}[title={What you've built}]
One more word for describing what you lived and felt.
\end{grammarlens}

\subsection*{Your turn}
\begin{itemize}
\raggedright
  \item \emph{Say it:} \emph{pyār}
  \item \emph{Say it:} \emph{pyār}, once more
  \item \emph{Say it:} \emph{mātānā pyārthī bāḷak śānt thaī gayũ.} — In its mother's love the child grew calm.
  \item \emph{Recall:} say the Gujarati for to kiss, then the Gujarati for to hold on tight, then say \emph{pyār} again
\end{itemize}

\begin{tcolorbox}[breakable,colback=teal!4,colframe=teal!35!black,title={Before you move on}]
What does \gu{પ્યાર} mean? (\textquotedblleft{}Love; loving fondness\textquotedblright{}.) Say it once more.
\end{tcolorbox}
\section[\texorpdfstring{\gu{ચાંપવું} (cā̃pvũ)}{ચાંપવું (cā̃pvũ)}]{\gu{ચાંપવું} (cā̃pvũ) — to press close; to clasp}

\label{lesson:GU-C451-capvu}



\begin{quote}
Before the new one: say the Gujarati for to hold on tight, then the Gujarati for love.
\end{quote}

\subsection*{You'll want to know: \gu{ચાંપવું}}
\textbf{\gu{ચાંપવું}} — \emph{cā̃pvũ} — \textquotedblleft{}to press close; to clasp\textquotedblright{}.

\begin{quote}
\gu{માતાએ} \gu{બાળકને} \gu{છાતીએ} \gu{ચાંપ્યું}.
\end{quote}

\begin{quote}
\emph{mātāe bāḷakne chātīe cā̃pyũ.}
\end{quote}

\begin{quote}
The mother clasped the child to her chest.
\end{quote}

Past, with \gu{તેણે} (he or she …): \gu{તેણે} \gu{બાળકને} \gu{છાતીએ} \gu{ચાંપ્યું}, \emph{teṇe bāḷakne chātīe cā̃pyũ}.

\begin{cousinweb}
\textbf{\gu{ચાંપવું}} (\emph{cā̃pvũ}) goes back to a reconstructed \emph{*capp-}, \textquotedblleft{}press\textquotedblright{}; CDIAL lists the Gujarati verb, \textquotedblleft{}to press\textquotedblright{} (CDIAL 4674).
\end{cousinweb}

\begin{grammarlens}[title={What you've built}]
One more word for describing what you lived and felt.
\end{grammarlens}

\subsection*{Your turn}
\begin{itemize}
\raggedright
  \item \emph{Say it:} \emph{cā̃pvũ}
  \item \emph{Say it:} \emph{cā̃pvũ}, once more
  \item \emph{Say it:} \emph{mātāe bāḷakne chātīe cā̃pyũ.} — The mother clasped the child to her chest.
  \item \emph{Recall:} say the Gujarati for to hold on tight, then the Gujarati for love, then say \emph{cā̃pvũ} again
\end{itemize}

\begin{tcolorbox}[breakable,colback=teal!4,colframe=teal!35!black,title={Before you move on}]
What does \gu{ચાંપવું} mean? (\textquotedblleft{}To press close; to clasp\textquotedblright{}.) Say it once more.
\end{tcolorbox}
